\section{Instantiation}\label{sec:tsps-instantiation}

This section instantiates the construction of \Cref{sec:tsps-construction}
twice, with $\AGHO$ and $\KPW$, giving full descriptions of two threshold SPS
schemes built from the MPC protocol~\cite{C:AGHO11,C:KilPanWee15}. In each case we verify that
the signing algorithm admits the decomposition the construction requires, and
give the resulting threshold scheme.

$\AGHO$ has the smallest possible signatures, and its secret key consists of
field elements only, so its threshold version only needs multiplication of
shared field elements. $\KPW$ relies on a standard assumption, but its secret
key contains group elements, so its threshold version needs sharings of group
elements and $\algo{ExpGrp}$. In both, the signers raise a public group element
to a shared exponent by each raising it to their own share, which gives a
sharing of the result.

\subsection{Instantiating the MPC protocol}\label{sec:tsps-mpc-inst}

We realise $\algo{ExpGrp}$ as Smart and Talibi Alaoui do, from multiplication
triples made in preprocessing~\cite{C:Beaver91b,IMA:SmaTal19}.

$\algo{ExpGrp}$ consumes a mixed triple $(\shr{\eta}, \shr{h}, \shr{h^\eta})$
with $\eta \sample \Zp$ and $h \sample \group{}$. The signers open the
following, which are uniform and reveal nothing:
%
\begin{equation*}
%
  \epsilon = \alpha - \eta, \quad \delta = b h^{-1}.
%
\end{equation*}
%
They then compute the following locally:
%
\begin{equation*}
%
  \shr{b^{\alpha}} = \shr{h^{\eta}} \cdot \shr{h}^{\epsilon} \cdot
\delta^{\shr{\eta}} \cdot \delta^{\epsilon},
%
\end{equation*}
%
where one designated signer adds the public term $\delta^{\epsilon}$. The
signers obtain a mixed triple from an ordinary triple over $\Zp$ by raising a
generator to the shares of its second and third entries. We added this to
MP-SPDZ, together with sharings of group elements and their authentication
information.

Neither protocol we benchmark is proven secure against adaptive corruption, so
\Cref{lem:tsps-adaptive} does not apply to them.

\subsection{Instantiating $\TSPS$ with $\AGHO$}\label{sec:tsps-abe}

Abe, Groth, Haralambiev, and Ohkubo established a lower bound of three group
elements for signature size in type III pairing groups, and gave a matching
construction that signs messages in both groups~\cite{C:AGHO11}. Its signatures
remain three group elements regardless of message length, and it is strongly
unforgeable in the GGM.

They also gave a rerandomisable three-element scheme for messages in one group
only, which is existentially unforgeable in the GGM~\cite{C:AGHO11}. We use this
scheme with $\group{1}$ and $\group{2}$ swapped, so that messages are in
$\group{1}$, and call it $\AGHO$. \Cref{prim:agho} recalls it. Its secret
key consists only of field elements, and the message enters signing only through
the terms $\msg_i^{-\mu_i}$. Our threshold version relies on both properties.

\begin{primitivebox}{The $\AGHO$ SPS scheme~\cite{C:AGHO11}.}{agho}

  \begin{description}

    \item[$(\sk, \pk) \sample \SPSgen(\secparam, \ell)$ :] Sample $\mu_1,
\ldots, \mu_\ell, \tau \sample \Zp^*$, and return the following:
      %
      \begin{gather*}
      %
	\sk = (\mu_1, \ldots, \mu_\ell, \tau), \quad \pk = (\widetilde{u}_1,
\ldots, \widetilde{u}_\ell, v), \\
      %
	\widetilde{u}_i = \ggen^{\mu_i} \ \text{ for } i \in [1,\ell], \quad v =
\gen^{\tau}.
      %
      \end{gather*}

    \item[$\SPSsig \sample \SPSsign(\sk, \vec{\msg})$ :] On input $\vec{\msg} =
(\msg_1, \ldots, \msg_\ell) \in \group{1}^\ell$, sample $\omega \sample \Zp^*$
and return the following:
      %
      \begin{equation*}
      %
	\SPSsig = (\widetilde{r}, \widetilde{s}, z) = \left(\ggen^{\omega}, \
\widetilde{r}^{\,\tau}, \ \left(\gen \prod_{i=1}^{\ell}
\msg_i^{-\mu_i}\right)^{1/\omega}\right).
      %
      \end{equation*}

    \item[$\SPSsig' \sample \SPSrandomise(\SPSsig)$ :] On input $\SPSsig =
(\widetilde{r}, \widetilde{s}, z)$, sample $\omega' \sample \Zp^*$ and return
$(\widetilde{r}^{\,\omega'}, \widetilde{s}^{\,\omega'}, z^{1/\omega'})$.

    \item[$\{0,1\} \leftarrow \SPSverify(\pk, \vec{\msg}, \SPSsig)$ :] On input
$\SPSsig = (\widetilde{r}, \widetilde{s}, z)$, return $1$ if and only if the
following hold:
      %
      \begin{equation*}
      %
	\pairing(v, \widetilde{r}) = \pairing(\gen, \widetilde{s}), \quad
\pairing(z, \widetilde{r}) \cdot \prod_{i=1}^{\ell} \pairing(\msg_i, \widetilde{u}_i)
= \pairing(\gen, \ggen).
      %
      \end{equation*}

  \end{description}

\end{primitivebox}

Writing $\rho =
1/\omega$ for the signing randomness of $\AGHO$, a signature on $\vec{\msg}$ is as
follows:
%
\begin{equation*}
%
  (\widetilde{r}, \widetilde{s}, z) = \left(\ggen^{1/\rho}, \
\ggen^{\tau/\rho}, \ \gen^{\rho} \prod_{i=1}^{\ell} \msg_i^{-\rho\mu_i}\right).
%
\end{equation*}
%
It has decomposable signing with $c_0$, $c_1$ and $e_{1,i}$ as follows:
%
\begin{equation*}
%
  c_0 = (\widetilde{r}, \widetilde{s}), \quad c_1 = \gen^{\rho}, \quad e_{1,i}
= -\rho\mu_i,
%
\end{equation*}
%
none of which depend on the message.
Everything the offline phase computes is a field element or a generator raised
to one, so it needs no $\algo{ExpGrp}$. The costly step is $1/\rho$. Computed
directly, it takes a multiplication and an opening, and $\tau/\rho$ then takes
one more multiplication, which makes three rounds. Multiplying $\tau$ by the same
mask as $\rho$ in the first round gives $\tau/\rho$ once the mask is opened, so
two rounds suffice. The result is threshold $\AGHO$, given in
\Cref{prot:tsps-abe}. Since $\rho$ and $\beta$ are both
sampled from $\Zp$, the opened $\gamma = \rho\beta$ is zero with probability
about $2/p$. Under adaptive corruption, the signers' shares of $\rho$, $\beta$
and $\tau\beta$ are among the intermediate shares they are assumed to erase.

\begin{protocolbox}{The threshold structure-preserving signature scheme from AGHO}{tsps-abe}

  \begin{description}

    \item[$(\shr{\sk}, \pk) \sample \TSPSgen(\secparam, \ell, t, n)$ :] The
signers do the following:

      \begin{enumerate}

	\item Sample $\shr{\mu_1}, \ldots, \shr{\mu_\ell}, \shr{\tau} \sample
\algo{RandShare}()$.

	\item Compute the following:
	  %
	  \begin{equation*}
	  %
	    \shr{v} = \gen^{\shr{\tau}}, \quad \shr{\widetilde{u}_i} = \ggen^{\shr{\mu_i}} \ \ (i \in [1,\ell]).
	  %
	  \end{equation*}

	\item Open $\pk = (\widetilde{u}_1, \ldots, \widetilde{u}_\ell, v)$.

	\item Each signer keeps $\shr{\sk} = (\shr{\mu_1}, \ldots,
\shr{\mu_\ell}, \shr{\tau})$.

      \end{enumerate}

    \item[$\shr{\state} \sample \TSPSoffsign(\shr{\sk})$ :] The signers do the
following:

      \begin{enumerate}

	\item Sample $\shr{\rho}, \shr{\beta} \sample \algo{RandShare}()$.

	\item Compute the following in one round:
	  %
	  \begin{gather*}
	  %
	    \shr{\rho\beta} \sample \algo{Mul}(\shr{\rho}, \shr{\beta}), \quad
\shr{\tau\beta} \sample \algo{Mul}(\shr{\tau}, \shr{\beta}), \\
	  %
	    \shr{\rho\mu_i} \sample \algo{Mul}(\shr{\rho}, \shr{\mu_i}) \ \ (i
\in [1,\ell]).
	  %
	  \end{gather*}

	\item In a second round, open $\gamma = \rho\beta$, and abort if $\gamma
= 0$.

	\item Each signer computes the following locally:
	  %
	  \begin{gather*}
	  %
	    \shr{1/\rho} = \gamma^{-1}\shr{\beta}, \quad \shr{\tau/\rho} =
\gamma^{-1}\shr{\tau\beta}, \\
	  %
	    \shr{\widetilde{r}} = \ggen^{\shr{1/\rho}}, \quad \shr{\widetilde{s}}
= \ggen^{\shr{\tau/\rho}}, \quad \shr{\gen^{\rho}} = \gen^{\shr{\rho}}.
	  %
	  \end{gather*}

	\item Store $\shr{\state} = (\shr{\widetilde{r}}, \shr{\widetilde{s}},
\shr{\gen^{\rho}}, \shr{e_{1,1}}, \ldots, \shr{e_{1,\ell}})$, where
$\shr{e_{1,i}} = -\shr{\rho\mu_i}$.

      \end{enumerate}

    \item[$\SPSsig_h \leftarrow \TSPSonsign(\sk_h, \state_h, \vec{\msg})$ :]
Each signer does the following:

      \begin{enumerate}

	\item Compute the following:
	  %
	  \begin{equation*}
	  %
	    \shr{z} = \shr{\gen^{\rho}} \cdot \prod_{i=1}^\ell
\msg_i^{\shr{e_{1,i}}}.
	  %
	  \end{equation*}

	\item Return $(\shr{\widetilde{r}}, \shr{\widetilde{s}}, \shr{z})$.

      \end{enumerate}

    \item[$\SPSsig \leftarrow \TSPScombine(\{\SPSsig_h\}_{h \in Q})$ :]
Reconstruct $\widetilde{r}$, $\widetilde{s}$ and $z$ from the shares of the $t+1$
signers in $Q$, and return $\SPSsig = (\widetilde{r}, \widetilde{s}, z)$.

    \item[$\{0,1\} \leftarrow \TSPSverify(\pk, \vec{\msg}, \SPSsig)$ :] Return
$1$ if and only if the following hold:
      %
      \begin{equation*}
      %
	\pairing(v, \widetilde{r}) = \pairing(\gen, \widetilde{s}), \quad
\pairing(z, \widetilde{r}) \cdot \prod_{i=1}^{\ell} \pairing(\msg_i, \widetilde{u}_i)
= \pairing(\gen, \ggen).
      %
      \end{equation*}

  \end{description}

\end{protocolbox}

\subsection{Instantiating $\TSPS$ with $\KPW$}\label{sec:tsps-kiltz}

\Cref{prim:sps-kiltz} recalls the $\KPW$ SPS scheme~\cite{C:KilPanWee15}. We write
$k_{i,j}$ for the entry of $\mat{K}$ in row $i$ and column $j$, and
$\mathsf{T}$ denotes transposition. We write $(\gen, \vec{\msg})$ for the
vector $(\gen, \msg_1, \ldots, \msg_\ell)$, and $\gen^{\vec{\theta}}$ for
$(\gen^{\theta_1}, \gen^{\theta_2}, \ldots)$ when $\vec{\theta}$ is a vector of
field elements. We take products of vectors entry by entry. For vectors
$\vec{f}$ and $\vec{\widetilde{f}}$ of the same length, with entries in
$\group{1}$ and $\group{2}$, we write the following:
%
\begin{equation*}
%
  \pairing(\vec{f}, \vec{\widetilde{f}}) = \prod_{i} \pairing(f_i,
\widetilde{f}_i).
%
\end{equation*}

\begin{primitivebox}{The $\KPW$ SPS scheme~\cite{C:KilPanWee15}.}{sps-kiltz}

  \begin{description}

    \item[$(\sk, \pk) \sample \SPSgen(\secparam, \ell)$ :] Sample $\alpha, \beta
\sample \Zp$, $\mat{K} \sample \Zp^{(\ell+1)\times 2}$ and $\mat{\Gamma},
\mat{\Delta} \sample \Zp^{2\times 2}$. Let $\vec{\alpha} = (1,
\alpha)^{\mathsf{T}}$ and $\vec{\beta} = (1, \beta)^{\mathsf{T}}$, and compute
the following:
      %
      \begin{gather*}
      %
	\vec{\mu} = \mat{K}\vec{\alpha}, \quad \vec{\nu} =
\mat{\Gamma}\vec{\alpha}, \quad \vec{\omega} = \mat{\Delta}\vec{\alpha}, \\
      %
	\vec{\xi} = \vec{\beta}^{\mathsf{T}}\mat{\Gamma}, \quad \vec{\upsilon} =
\vec{\beta}^{\mathsf{T}}\mat{\Delta}.
      %
      \end{gather*}
      %
      Return the following:
      %
      \begin{gather*}
      %
	\sk = (\mat{K}, \vec{b}, \vec{x}, \vec{y}) = (\mat{K}, \gen^{(1,
\beta)}, \gen^{\vec{\xi}}, \gen^{\vec{\upsilon}}), \\
      %
	\pk = (\vec{\widetilde{a}}, \vec{\widetilde{u}}, \vec{\widetilde{v}},
\vec{\widetilde{w}}) = (\ggen^{(1, \alpha)}, \ggen^{\vec{\mu}},
\ggen^{\vec{\nu}}, \ggen^{\vec{\omega}}).
      %
      \end{gather*}

    \item[$\SPSsig \sample \SPSsign(\sk, \vec{\msg})$ :] Sample $\rho, \sigma
\sample \Zp^*$ and return $\SPSsig = (\widetilde{q}, \vec{r}, \vec{s},
\vec{z})$, computed as follows:
      %
      \begin{equation*}
      %
	\widetilde{q} = \ggen^{\sigma}, \quad \vec{r} = \vec{b}^{\,\rho},
\quad \vec{s} = \vec{b}^{\,\rho\sigma}, \quad z_j = \gen^{k_{1,j}}
\prod_{i=1}^{\ell} \msg_i^{k_{i+1,j}} \cdot x_j^{\rho}\, y_j^{\rho\sigma}
\text{ for } j \in \{1, 2\}.
      %
      \end{equation*}

    \item[$\{0,1\} \leftarrow \SPSverify(\pk, \vec{\msg}, \SPSsig)$ :] Return
$1$ if and only if the following hold:
      %
      \begin{equation*}
      %
	\pairing(\vec{z}, \vec{\widetilde{a}}) = \pairing((\gen,
\vec{\msg}), \vec{\widetilde{u}}) \cdot \pairing(\vec{r},\vec{\widetilde{v}}) \cdot
\pairing(\vec{s}, \vec{\widetilde{w}}), \quad \pairing(r_1, \widetilde{q})
= \pairing(s_1, \ggen),
      %
      \end{equation*}
      %
      where $r_1$ and $s_1$ are the first entries of $\vec{r}$ and $\vec{s}$.

  \end{description}

\end{primitivebox}

$\KPW$ is EUF-CMA secure under SXDH~\cite{C:KilPanWee15}. It has decomposable
signing with $c_0$ and, for entry $j$ of $\vec{z}$, $c_j$ and $e_{j,i}$ as
follows:
%
\begin{equation*}
%
  c_0 = (\widetilde{q}, \vec{r}, \vec{s}), \quad c_j = \gen^{k_{1,j}}
x_j^{\rho} y_j^{\rho\sigma}, \quad e_{j,i} = k_{i+1,j}.
%
\end{equation*}
%
None of these depend on the message. Our protocol and implementation
compute $\gen^{k_{1,j}}$ in the online phase instead, which is also local.

The group elements of the secret key are $\vec{b}, \vec{x}, \vec{y} \in
\group{1}^2$, and the offline phase raises each of them to a shared exponent
with $\algo{ExpGrp}_1$.

\Cref{prot:tsps-kiltz} gives threshold $\KPW$. Its offline phase makes one
multiplication, then one round of eight calls to $\algo{ExpGrp}_1$.

\begin{protocolbox}{The threshold structure-preserving signature scheme from KPW}{tsps-kiltz}

  \begin{description}

    \item[$(\shr{\sk}, \pk) \sample \TSPSgen(\secparam, \ell, t, n)$ :] The
signers do the following:

      \begin{enumerate}

	\item Sample $\shr{\alpha}$, $\shr{\beta}$ and every entry of
$\shr{\mat{K}}$, $\shr{\mat{\Gamma}}$ and $\shr{\mat{\Delta}}$ with
$\algo{RandShare}$.

	\item Compute $\shr{\vec{\mu}}$, $\shr{\vec{\nu}}$, $\shr{\vec{\omega}}$,
$\shr{\vec{\xi}}$ and $\shr{\vec{\upsilon}}$ with one call to $\algo{Mul}$ per
entry, followed by a local addition, since the first entries of $\vec{\alpha}$ and
$\vec{\beta}$ are the public constant $1$.

	\item Each signer computes the following locally:
	  %
	  \begin{gather*}
	  %
	    \shr{\vec{b}} = \gen^{(1, \shr{\beta})}, \quad \shr{\vec{x}} =
\gen^{\shr{\vec{\xi}}}, \quad \shr{\vec{y}} = \gen^{\shr{\vec{\upsilon}}}, \\
	  %
	    \shr{\vec{\widetilde{a}}} = \ggen^{(1, \shr{\alpha})}, \quad
\shr{\vec{\widetilde{u}}} = \ggen^{\shr{\vec{\mu}}}, \quad
\shr{\vec{\widetilde{v}}} = \ggen^{\shr{\vec{\nu}}}, \quad
\shr{\vec{\widetilde{w}}} = \ggen^{\shr{\vec{\omega}}}.
	  %
	  \end{gather*}

	\item Open $\pk = (\vec{\widetilde{a}}, \vec{\widetilde{u}},
\vec{\widetilde{v}}, \vec{\widetilde{w}})$.

	\item Each signer keeps $\shr{\sk} = (\shr{\mat{K}}, \shr{\vec{b}},
\shr{\vec{x}}, \shr{\vec{y}})$.

      \end{enumerate}

    \item[$\shr{\state} \sample \TSPSoffsign(\shr{\sk})$ :] The signers do the
following:

      \begin{enumerate}

	\item Sample $\shr{\rho}, \shr{\sigma} \sample \algo{RandShare}()$.

	\item Compute the following:
	  %
	  \begin{gather*}
	  %
	    \shr{\rho\sigma} \sample \algo{Mul}(\shr{\rho}, \shr{\sigma}), \\
	  %
	    \shr{\vec{r}} \sample \algo{ExpGrp}_1(\shr{\rho}, \shr{\vec{b}}),
\quad \shr{\vec{s}} \sample \algo{ExpGrp}_1(\shr{\rho\sigma}, \shr{\vec{b}}),
\\
	  %
	    \shr{\vec{x}^{\,\rho}} \sample \algo{ExpGrp}_1(\shr{\rho},
\shr{\vec{x}}), \quad \shr{\vec{y}^{\,\rho\sigma}} \sample
\algo{ExpGrp}_1(\shr{\rho\sigma}, \shr{\vec{y}}).
	  %
	  \end{gather*}

	\item Each signer computes the following locally:
	  %
	  \begin{equation*}
	  %
	    \shr{\widetilde{q}} = \ggen^{\shr{\sigma}}, \quad \shr{\vec{z}^*} =
\shr{\vec{x}^{\,\rho}} \cdot \shr{\vec{y}^{\,\rho\sigma}}.
	  %
	  \end{equation*}

	\item Store $\shr{\state} = (\shr{\widetilde{q}}, \shr{\vec{r}},
\shr{\vec{s}}, \shr{\vec{z}^*})$.

      \end{enumerate}

    \item[$\SPSsig_h \leftarrow \TSPSonsign(\sk_h, \state_h, \vec{\msg})$ :]
Each signer does the following:

      \begin{enumerate}

	\item Compute the following:
	  %
	  \begin{equation*}
	  %
	    \shr{z_j} = \shr{z^*_j} \cdot \gen^{\shr{k_{1,j}}} \prod_{i=1}^{\ell}
\msg_i^{\shr{k_{i+1,j}}} \ \text{ for } j \in \{1, 2\}.
	  %
	  \end{equation*}

	\item Return $(\shr{\widetilde{q}}, \shr{\vec{r}}, \shr{\vec{s}},
\shr{\vec{z}})$.

      \end{enumerate}

    \item[$\SPSsig \leftarrow \TSPScombine(\{\SPSsig_h\}_{h \in Q})$ :]
Reconstruct every group element from the shares of the $t+1$ signers in $Q$, and
return $\SPSsig = (\widetilde{q}, \vec{r}, \vec{s}, \vec{z})$.

    \item[$\{0,1\} \leftarrow \TSPSverify(\pk, \vec{\msg}, \SPSsig)$ :] Return
$1$ if and only if the following hold:
      %
      \begin{equation*}
      %
	\pairing(\vec{z}, \vec{\widetilde{a}}) = \pairing((\gen,
\vec{\msg}), \vec{\widetilde{u}}) \cdot \pairing(\vec{r},\vec{\widetilde{v}}) \cdot
\pairing(\vec{s}, \vec{\widetilde{w}}), \quad \pairing(r_1, \widetilde{q})
= \pairing(s_1, \ggen).
      %
      \end{equation*}

  \end{description}

\end{protocolbox}

\begin{theorem}\label{thm:tsps-security} Let the MPC protocol be private and
secure with abort against up to $t$ statically corrupted signers, over a linear
$(t+1, n)$ threshold sharing. Then
the threshold $\AGHO$ and $\KPW$ schemes are correct. The threshold $\AGHO$
scheme is TS-UF-1 secure in the GGM, and the threshold $\KPW$ scheme is
TS-UF-1 secure under SXDH, in both cases with no loss. If the MPC protocol is
secure against adaptive corruption, both are adp-TS-UF-1 secure with the loss of
\Cref{lem:tsps-adaptive}.\end{theorem}

We prove \Cref{thm:tsps-security} in \Cref{sec:tsps-proof-instantiation}.

